\chapter{Your First \LaTeX{} Document}
\label{chap:first-document}
\index{first document}\index{TeX distribution}\index{Overleaf}\index{compilation}

The best first document is not a template with forty lines of unexplained
settings. It is a page that you understand completely. We shall create that
page now. Choose the online path if you want to begin without installing
software, or the local path if a TeX system is already available. Both paths
lead to the same \file{main.tex} source and the same essential PDF.

\begin{learningoutcomes}
  \item Create a new \LaTeX{} project online or on a local computer.
  \item Identify the project's main source file.
  \item Write the smallest useful document with the \code{article} class.
  \item Compile the source and inspect the PDF preview.
  \item Explain each line of the minimal document.
  \item Locate the log when compilation fails.
  \item Save or download both source and PDF output.
\end{learningoutcomes}

\section{Choose a workspace}

A \term{\LaTeX{} project} is the group of files needed to build a document. Our
first project needs only one source file. Later projects will also contain
figures, bibliography data, and separate section files.

\subsection{The online path: Overleaf}

Overleaf runs the \LaTeX{} tools on a remote service and displays an editor and
PDF preview in a web browser. It is convenient for the first lesson because
there is no local installation to diagnose. The names and positions of
interface controls may change, so treat the steps below as tasks rather than a
permanent map of the screen. The current official documentation explains
project creation, recompilation, main documents, and downloads
\parencite{overleafFirstProject,overleafRecompile,overleafMainDocument,
overleafDownload}.

\begin{enumerate}
  \item Sign in to Overleaf and create a new blank project.
  \item Give it a useful name, such as \file{first-latex-document}.
  \item In the file list, select \file{main.tex}. Delete any sample content so
        that the file is empty.
  \item Type the source shown in \cref{sec:minimal-document} exactly.
  \item Run the project's Recompile action.
  \item Inspect the PDF preview and, if necessary, open the log or error panel.
\end{enumerate}

Overleaf treats one \file{.tex} file as the \term{main document}: the file sent
to the compiler first. In a one-file project, \file{main.tex} is the natural
choice. In a multi-file project, a chapter file is not normally the main
document because it may depend on settings stored elsewhere.

\begin{classroomnote}
Overleaf normally saves edits automatically, but a remotely saved project is
not a substitute for an independent backup. Download the source as a ZIP file
at sensible milestones and download the PDF separately. Current Overleaf
documentation notes that the ordinary source download and PDF download are
separate actions \parencite{overleafDownload}.
\end{classroomnote}

\subsection{The local-computer path}

Local work requires a \term{TeX distribution}: a maintained collection of the
compiler, standard classes, packages, fonts, and utilities. Common choices are
TeX Live on macOS and Linux, and either MiKTeX or TeX Live on Windows. macOS
users often install MacTeX, which is a TeX Live distribution packaged for the
platform. Installation details change, so use the current official installer
for your operating system rather than an old third-party tutorial.

You also need a plain-text editor. TeXstudio is a focused \LaTeX{} editor. Visual
Studio Code can become a \LaTeX{} editor through the \LaTeX{} Workshop extension.
Neither editor changes the language in \file{main.tex}; it changes how you
edit, compile, and view it.

After installing a distribution and editor:

\begin{enumerate}
  \item Create a folder named \file{first-latex-document}.
  \item In that folder, create a plain-text file named \file{main.tex}.
  \item Type the minimal source in \cref{sec:minimal-document}.
  \item Save the file before compiling.
  \item Use the editor's build action, or run \pkg{latexmk} as described in
        \cref{sec:latexmk-first}.
\end{enumerate}

Do not write the source in an application that silently saves a word-processing
format. The file must be plain text and its name must end in \file{.tex}, not
\file{.tex.txt}.

\begin{goodpractice}
Use a short folder name made from letters, numbers, hyphens, or underscores.
For example, \file{first-latex-document} travels more reliably than a name
containing punctuation and several spaces. Later we will develop a complete
file-naming policy for collaborative projects.
\end{goodpractice}

\section{The smallest useful document}
\label{sec:minimal-document}

Type these four lines. Do not add a title, packages, page settings, or copied
preamble material yet.

\begin{latexcode}{Complete example: main.tex}
\documentclass{article}
\begin{document}
Hello, \LaTeX{}!
\end{document}
\end{latexcode}

Save the source and compile it. A successful build produces a PDF. If your
editor shows a green tick or a message such as ``process exited normally'',
that status is helpful, but the PDF itself still deserves inspection.

\expectedoutput

You should see one page containing ``Hello, \LaTeX{}!'' as an ordinary line of
text. The type, margins, page size, and page number come from class defaults.
Your result may differ slightly across current TeX installations, but the
sentence should be present and readable.

\sourceanatomy

\paragraph{Line 1: \cmd{documentclass}\required{article}.}
The backslash begins a command. The command name is
\code{documentclass}. The braces contain a required piece of information: the
class name \code{article}. A class acts like an architectural plan. It defines
the broad structures and default appearance available to the document.

\paragraph{Line 2: \cmd{begin}\required{document}.}
This line opens the \code{document} \term{environment}. An environment is a
region with a defined beginning and end. Here it marks the beginning of
material that can appear in the output.

\paragraph{Line 3: ordinary text.}
This line contains no command. \LaTeX{} treats it as text to typeset. We used a
comma and exclamation mark in the ordinary way.

\paragraph{Line 4: \cmd{end}\required{document}.}
This line closes the environment opened on line 2. \LaTeX{} ignores source after
this line during an ordinary document build, so the document should genuinely
end here.

The portion before \cmd{begin}\required{document} is called the
\term{preamble}. The portion between the opening and closing lines is the
\term{document body}. Our preamble currently contains only the class
declaration.

\begin{tryit}
Replace the sentence with your own plain sentence, save the file, and compile
again. Confirm that the sentence changes but the page's broad design does not.
Then add a second sentence on the next source line and compile once more. What
does \LaTeX{} do with the line break?
\end{tryit}

The second source line normally continues the same paragraph. A source line is
not automatically a printed line. This behaviour lets us wrap source for
readability without commanding a visual break after every line. Blank lines
and paragraphs receive a full explanation in \cref{chap:source-anatomy}.

\section{What compilation does}

When you compile, the engine reads the class declaration first. It then reads
the body in order, interprets commands, and constructs pages. It writes a PDF
and several temporary \term{auxiliary files}. An auxiliary file stores
information needed by the build, such as a table-of-contents entry or a label
number.

This explains why a longer document sometimes needs more than one compiler
pass. On the first pass, \LaTeX{} records information. On a later pass, it reads
that information and places it in the PDF. You do not need to manage those
passes by hand when a tool such as \pkg{latexmk} can detect them.

\subsection{A first local build with latexmk}
\label{sec:latexmk-first}

A \term{terminal} is a text-based window in which you give the computer a
command. Open a terminal in the folder containing \file{main.tex}, then enter:

\begin{terminalcode}{Terminal command}
latexmk -pdf main.tex
\end{terminalcode}

Here \code{latexmk} is the build tool, \code{-pdf} requests PDF output, and
\file{main.tex} is the main source. \pkg{latexmk} runs the required programs
again when it detects unfinished references or bibliography work. Its official
manual describes this automated build role \parencite{latexmkCTAN}.

To remove generated auxiliary files after an experiment, use:

\begin{terminalcode}{Clean generated files}
latexmk -C main.tex
\end{terminalcode}

The capital \code{-C} removes generated products recognised by
\pkg{latexmk}. Check that you are in the intended project folder before using
any clean command. Your \file{main.tex} source should remain.

\begin{advancedtopic}
An editor's build button usually runs one or more command-line programs behind
the interface. Learning the \pkg{latexmk} command is useful because it gives a
project a reproducible instruction that does not depend on one editor. You may
continue using the build button while you are learning.
\end{advancedtopic}

\section{Compile, inspect, revise}

Compilation succeeding is necessary, but it is not the same as the document
being correct. Develop a four-part routine:

\begin{enumerate}
  \item \textbf{Compile} after a small, coherent edit.
  \item \textbf{Inspect} the area that should have changed.
  \item \textbf{Read} any error and the most relevant warnings.
  \item \textbf{Revise} the source, not the generated PDF.
\end{enumerate}

Frequent builds keep cause and effect close together. If you write 200 lines
before compiling, an error could be anywhere in those lines. If you write five
lines, the search area is much smaller.

\begin{pausepredict}
Change \code{article} to \code{book}, but do not compile yet. Will the words
``Hello, \LaTeX{}!'' remain? Might some broad page choices change? Make a
prediction, compile, and compare.

Then restore \code{article}. We will compare classes properly after learning
the structures they provide.
\end{pausepredict}

\section{Your first controlled error}

Keep a working copy of the successful source. In the practice copy, remove
the backslash from the second line:

\begin{latexcode}{Broken example: do not use as a working document}
\documentclass{article}
begin{document}
Hello, \LaTeX{}!
\end{document}
\end{latexcode}

Compile. The wording varies with the engine and editor, but the log should
report a problem involving a missing \cmd{begin}\required{document} or text in
the wrong place. \LaTeX{} saw the characters \code{begin\{document\}} as text,
not as the command required to open the body.

Repair the line:

\begin{latexcode}{Corrected line}
\begin{document}
\end{latexcode}

Compile again and confirm that the successful PDF returns.

\begin{commonerror}
\textbf{Symptom:} the editor reports ``Missing \cmd{begin}\required{document}''.

\textbf{Likely causes:} the opening command is missing, misspelt, or lacks its
backslash; or visible text has accidentally been placed in the preamble.

\textbf{Repair:} compare the top of the file with the four-line working
example. Correct the first structural difference, save, and compile again.
Do not respond by adding several unexplained commands at once.
\end{commonerror}

\section{Saving and keeping both forms}

The editable scholarly record is not just the PDF. Keep \file{main.tex} and,
when the project grows, all figures, data used to generate figures, bibliography
files, local class/style files permitted for sharing, and build instructions.
The PDF records what a particular build produced; the source makes revision
and reproduction possible.

On Overleaf, use the current project menu to download source and the PDF-view
control to download the typeset PDF. The source download is normally a ZIP
archive. Extract it and verify that \file{main.tex} is present. On a local
computer, make a backup outside the project folder or in a properly managed
versioned location.

\begin{scienceinpractice}
Suppose an environmental monitoring report is revised after one unit label is
corrected. Keeping only the old PDF preserves the error. Keeping source plus
the data, figure, and build instruction makes the correction traceable and the
revised report reproducible.
\end{scienceinpractice}

\section{The first complete working pattern}

We can now state our first eight-step pattern explicitly:

\begin{enumerate}
  \item \textbf{Identify the problem:} create one typeset page.
  \item \textbf{Describe the output:} one sentence in a readable PDF.
  \item \textbf{Select the structure:} the standard \code{article} class and
        document environment.
  \item \textbf{Write the smallest source:} four lines.
  \item \textbf{Compile and inspect:} confirm the sentence appears.
  \item \textbf{Correct errors:} compare the first structural difference.
  \item \textbf{Improve carefully:} change one sentence, then rebuild.
  \item \textbf{Check fitness:} retain both source and PDF. This is a practice
        page, not yet a submission-ready article.
\end{enumerate}

That last sentence matters. A file can compile successfully without meeting a
scholarly purpose. Technical validity and publication readiness are related,
but they are not identical.

\chapterproject

\begin{miniproject}
Inside \file{measurement-methods-article}, create \file{main.tex} with this
source:

\begin{latexcode}{Running project: first compilable version}
\documentclass{article}
\begin{document}
This project compares Method A and Method B.

All values are simulated demonstration data.
\end{document}
\end{latexcode}

Compile it. The blank source line begins a new paragraph. Inspect the PDF and
keep both source and output. Add a plain-text \file{README.md} stating that the
project is educational and the data are simulated. Do not add titles,
sections, equations, or packages yet; we shall earn each layer.
\end{miniproject}

\chaptersummary

A first project can be created online or locally. Its main source file is the
file the compiler reads first. The minimal document selects the
\code{article} class, opens the document body, contains ordinary text, and
closes the body. Compilation produces a PDF and auxiliary files. Frequent
small builds make errors easier to locate, and \pkg{latexmk} automates repeated
passes. A responsible archive keeps editable source as well as output.

\chaptercheck

\begin{chapterchecklist}
  \item I have compiled the four-line document successfully.
  \item I can point to the preamble and document body.
  \item I can explain every line without calling it ``just required code''.
  \item I know where my editor or online project shows the build log.
  \item I deliberately created and repaired the missing-backslash error.
  \item I saved or downloaded both source and PDF.
  \item My running project clearly labels its content as simulated.
\end{chapterchecklist}

\chapterexercisestitle

\begin{chapterexercises}
  \item Write the four-line minimal document from memory, then compare it with
        the printed source before compiling.
  \item What is the difference between the \code{article} class and the
        \code{document} environment in this example?
  \item Predict what happens when two ordinary source lines have no blank line
        between them. Test your prediction.
  \item Change ``Hello, \LaTeX{}!'' to a sentence from your field. Compile and
        identify which source line changed the output.
  \item Repair this broken opening without adding new commands:
        \code{begin\{document\}}.
  \item Explain why a successful compiler exit does not prove that a report is
        scientifically or typographically correct.
  \item Challenge: compile the same source online and locally, if both paths
        are available. Record any visible difference and avoid assuming that a
        difference is an error.
\end{chapterexercises}
